%----------------------------------------------------------------------------------------
%	PACKAGES AND OTHER DOCUMENT CONFIGURATIONS
%----------------------------------------------------------------------------------------
\documentclass[12pt]{article}
%----------------------------------------------------------------------------------------
%	PACKAGES AND OTHER DOCUMENT CONFIGURATIONS
%----------------------------------------------------------------------------------------

\usepackage{amsmath,amsfonts,stmaryrd,amssymb} % Math packages
\usepackage{enumerate} % Custom item numbers for enumerations
\usepackage[ruled]{algorithm2e} % Algorithms
\usepackage[framemethod=tikz]{mdframed} % Allows defining custom boxed/framed environments
\usepackage{textcomp}
\usepackage{float}
\usepackage{biblatex}
\usepackage{listings} % File listings, with syntax highlighting
\usepackage{graphicx}
\usepackage{hyperref}
\usepackage{ragged2e}
\usepackage{helvet}
\usepackage{rotating}
\usepackage[normalem]{}
\usepackage{colortbl}
\usepackage{multirow}
\usepackage{longtable}

\newcommand\ackname{Acknowledgements}
\if@titlepage
   \newenvironment{acknowledgements}{%
       \titlepage
       \null\vfil
       \@beginparpenalty\@lowpenalty
       \begin{center}%
         \bfseries \ackname
         \@endparpenalty\@M
       \end{center}}%
      {\par\vfil\null\endtitlepage}
\else
   \newenvironment{acknowledgements}{%
       \if@twocolumn
         \section*{\abstractname}%
       \else
         \small
         \begin{center}%
           {\bfseries \ackname\vspace{-.5em}\vspace{\z@}}%
         \end{center}%
         \quotation
       \fi}
       {\if@twocolumn\else\endquotation\fi}
\fi
\makeatother
%----------------------------------------------------------------------------------------
%	DOCUMENT MARGINS
%----------------------------------------------------------------------------------------

\usepackage{geometry} % Required for adjusting page dimensions and margins

\geometry{
	paper=a4paper, % Paper size, change to letterpaper for US letter size
	top=2.5cm, % Top margin
	bottom=3cm, % Bottom margin
	left=2.5cm, % Left margin
	right=2.5cm, % Right margin
	headheight=14pt, % Header height
	footskip=1.5cm, % Space from the bottom margin to the baseline of the footer
	headsep=1.2cm, % Space from the top margin to the baseline of the header
	%showframe, % Uncomment to show how the type block is set on the page
}

%----------------------------------------------------------------------------------------
%	FONTS
%----------------------------------------------------------------------------------------

\usepackage[utf8]{inputenc} % Required for inputting international characters
\usepackage[T1]{fontenc} % Output font encoding for international characters

\usepackage{XCharter} % Use the XCharter fonts

%----------------------------------------------------------------------------------------
%	COMMAND LINE ENVIRONMENT
%----------------------------------------------------------------------------------------

% Usage:
% \begin{commandline}
%	\begin{verbatim}
%		$ ls
%		
%		Applications	Desktop	...
%	\end{verbatim}
% \end{commandline}

\mdfdefinestyle{commandline}{
	leftmargin=10pt,
	rightmargin=10pt,
	innerleftmargin=15pt,
	middlelinecolor=black!50!white,
	middlelinewidth=2pt,
	frametitlerule=false,
	backgroundcolor=black!5!white,
	frametitle={Command Line},
	frametitlefont={\normalfont\sffamily\color{white}\hspace{-1em}},
	frametitlebackgroundcolor=black!50!white,
	nobreak,
}

% Define a custom environment for command-line snapshots
\newenvironment{commandline}{
	\medskip
	\begin{mdframed}[style=commandline]
}{
	\end{mdframed}
	\medskip
}

%----------------------------------------------------------------------------------------
%	FILE CONTENTS ENVIRONMENT
%----------------------------------------------------------------------------------------

% Usage:
% \begin{file}[optional filename, defaults to "File"]
%	File contents, for example, with a listings environment
% \end{file}

\mdfdefinestyle{file}{
	innertopmargin=1.6\baselineskip,
	innerbottommargin=0.8\baselineskip,
	topline=false, bottomline=false,
	leftline=false, rightline=false,
	leftmargin=2cm,
	rightmargin=2cm,
	singleextra={%
		\draw[fill=black!10!white](P)++(0,-1.2em)rectangle(P-|O);
		\node[anchor=north west]
		at(P-|O){\ttfamily\mdfilename};
		%
		\def\l{3em}
		\draw(O-|P)++(-\l,0)--++(\l,\l)--(P)--(P-|O)--(O)--cycle;
		\draw(O-|P)++(-\l,0)--++(0,\l)--++(\l,0);
	},
	nobreak,
}

% Define a custom environment for file contents
\newenvironment{file}[1][File]{ % Set the default filename to "File"
	\medskip
	\newcommand{\mdfilename}{#1}
	\begin{mdframed}[style=file]
}{
	\end{mdframed}
	\medskip
}

%----------------------------------------------------------------------------------------
%	NUMBERED QUESTIONS ENVIRONMENT
%----------------------------------------------------------------------------------------

% Usage:
% \begin{question}[optional title]
%	Question contents
% \end{question}

\mdfdefinestyle{question}{
	innertopmargin=1.2\baselineskip,
	innerbottommargin=0.8\baselineskip,
	roundcorner=5pt,
	nobreak,
	singleextra={%
		\draw(P-|O)node[xshift=1em,anchor=west,fill=white,draw,rounded corners=5pt]{%
		Question \theQuestion\questionTitle};
	},
}

\newcounter{Question} % Stores the current question number that gets iterated with each new question

% Define a custom environment for numbered questions
\newenvironment{question}[1][\unskip]{
	\bigskip
	\stepcounter{Question}
	\newcommand{\questionTitle}{~#1}
	\begin{mdframed}[style=question]
}{
	\end{mdframed}
	\medskip
}

%----------------------------------------------------------------------------------------
%	WARNING TEXT ENVIRONMENT
%----------------------------------------------------------------------------------------

% Usage:
% \begin{warn}[optional title, defaults to "Warning:"]
%	Contents
% \end{warn}

\mdfdefinestyle{warning}{
	topline=false, bottomline=false,
	leftline=false, rightline=false,
	nobreak,
	singleextra={%
		\draw(P-|O)++(-0.5em,0)node(tmp1){};
		\draw(P-|O)++(0.5em,0)node(tmp2){};
		\fill[black,rotate around={45:(P-|O)}](tmp1)rectangle(tmp2);
		\node at(P-|O){\color{white}\scriptsize\bf !};
		\draw[very thick](P-|O)++(0,-1em)--(O);%--(O-|P);
	}
}

% Define a custom environment for warning text
\newenvironment{warn}[1][Warning:]{ % Set the default warning to "Warning:"
	\medskip
	\begin{mdframed}[style=warning]
		\noindent{\textbf{#1}}
}{
	\end{mdframed}
}

%----------------------------------------------------------------------------------------
%	INFORMATION ENVIRONMENT
%----------------------------------------------------------------------------------------

% Usage:
% \begin{info}[optional title, defaults to "Info:"]
% 	contents
% 	\end{info}

\mdfdefinestyle{info}{%
	topline=false, bottomline=false,
	leftline=false, rightline=false,
	nobreak,
	singleextra={%
		\fill[black](P-|O)circle[radius=0.4em];
		\node at(P-|O){\color{white}\scriptsize\bf i};
		\draw[very thick](P-|O)++(0,-0.8em)--(O);%--(O-|P);
	}
}

% Define a custom environment for information
\newenvironment{info}[1][Info:]{ % Set the default title to "Info:"
	\medskip
	\begin{mdframed}[style=info]
		\noindent{\textbf{#1}}
}{
	\end{mdframed}
}
\bibliography{biblio}
\begin{document}
%----------------------------------------------------------------------------------------
%	TITLE PAGE
%----------------------------------------------------------------------------------------
\begin{titlepage}
	\newcommand{\HRule}{\rule{\linewidth}{0.5mm}}
	\center
	%------------------------------------------------
	%	Headings
	%------------------------------------------------
	\textsc{\LARGE Université de Namur}\\[1.5cm]
	\textsc{\Large Sécurité et Fiabilité des Systèmes Informatiques }\\[0.5cm]
	\textsc{\large IHDCM035}\\[0.5cm]
	%------------------------------------------------
	%	Title
	%------------------------------------------------	
	\HRule\\[0.4cm]
	{\huge\bfseries Etude des Risques: Informatisation d'un Centre Hospitalier}\\[0.4cm]
	\HRule\\[1.5cm]	
	%------------------------------------------------
	%	Author(s)
	%------------------------------------------------	
	\begin{minipage}{0.4\textwidth}
		\begin{flushleft}
			\large
			\textit{Auteur}\\
			Kenny \textsc{Warszawski} 
		\end{flushleft}
	\end{minipage}
	~
	\begin{minipage}{0.4\textwidth}
		\begin{flushright}
			\large
			\textit{Professeur}\\
			Jean-Nöel \textsc{Colin} 
		\end{flushright}
	\end{minipage}	
	%------------------------------------------------
	%	Date
	%------------------------------------------------	
	\vfill\vfill\vfill
	{\large\today}
	%------------------------------------------------
	%	Logo
	%------------------------------------------------
	\vfill\vfill
	\includegraphics[width=0.2\textwidth]{assets/placeholder.png}\\[1cm]
	\vfill
\end{titlepage}

\newpage
%----------------------------------------------------------------------------------------
%	Content
%----------------------------------------------------------------------------------------
\renewcommand{\contentsname}{Table des matières}
\tableofcontents
\newpage
%----------------------------------------------------------------------------------------
%	Introduction
%----------------------------------------------------------------------------------------

\section{Introduction} 

\subsection{Contexte}

\justify
Cette étude des risques concerne le Centre Hospitalier Mercy West(CHMW). Ce centre a mis en place un système informatique qui permet de centraliser les données de leurs patients. Afin de réaliser cela, l'hopital a mis à disposition un ordinateur connecté à une plateforme en ligne. Ainsi, le corps médical peut encoder les informations nécessaires sur leurs patients à la fin de leur service. Avant de commencer leur journée, le personnel peut également accéder aux dernières informations récoltées par leurs collègues pour rester à jour sur: l'état de santé des patients, les soins reçus, les opérations subies, les médicaments prescris, etc.

\justify
Chaque membre du personnel possède un badge afin de s'authentifier sur la plateforme. Les droits de lecture et modification d'un dossier médical sont associés à des droits qui sont assignés aux utilisateurs.  Ces droits sont associés à la fonction professionelle que l'utilisateur authentifié exerce. Par exemple, si un médecin s'authentifie, il pourra modifier les prescriptions de médicaments d'un patient tandis qu'une aide soignante ne pourra pas. Par contre, cette dernière aura le droit de modifier l'état de santé général du patient: taille, poids, nutrition, etc.

\justify
Ce logiciel impacte donc le quotidien des employés de cet hopital. Il est indispensable que tout le personnel indique rigoureusement les information concernant le patient. Ainsi, il sera possible de garantir un suivi médical journalier de haute qualité mais également d'en conserver un historique. Si une personne mal intentionnée ou tout simplement non-formée indique des informations erronées, ces peuvent avoir de lourdes conséquences sur le patient. Par conséquent, la réputation de l'hopital pourrait gravement diminuer. L'accès aux informations médicales, l'encodage des données ainsi que la réputation de l'hopital sont des \textbf{biens essentiels}. 

\justify
La confidentialité est un des critères de sécurité les plus important pour l'hôpital. De fait, si les informations médicales d'un patient arrivent entre de mauvaises mains, cela peut également avoir des conséquences dramatiques. Il est essentiel que les données médicales soient sécurisées et exploitable uniquement par les utilisateurs qui en ont le droit.

\justify
En ce qui concerne les \textbf{biens supports}, le centre hospitalier possède une infrastructure informatique dédiée afin de faire fonctionner l'ensemble de ses logiciels. Cette infrastructure comprend: des ordinateurs, des serveurs, un sous-réseau, un système d'authentification et de multiples disques durs afin de pouvoir stocker les données.

\subsection{Objectifs}

\justify
L'objectif de cette étude est de pouvoir établir une analyse de risque concernant ce projet. De plus, un plan d'action sera proposé en réponse aux évènements redoutés par le centre hospitalier. Le champs de cette étude sera toutefois limitée uniquement à la plateforme précédemment mentionnée. Tous les autres processus organisationnel ou informatiques nullement liés à la problématique ne seront pas pris en compte.

\section{Analyse de risques}

\subsection{Evènement redoutés}

Cette section est dédiée à une analyse des évènements redoutés. Cette analyse est basée sur les biens essentiels de l'hôpital et des critères de sécurités importants. (Disponibilité, Confidentialité et Intégrité)

\subsubsection{Accès aux informations médicales}

L'analyse de ce bien essentiel concerne la consultation des informations des patients. Par exemple, en début de service par un membre du corps médical.

\begin{table}[ht]
\centering
\resizebox{\textwidth}{!}{
\begin{tabular}{|l|l|l|l|l|}
	\hline
	Evenements Redoutés & Critère de Sécurité & Source de la Menace & Impact  & Severité \\ \hline

	\begin{tabular}[c]{@{}l@{}}Incendie dans la salle des\\ serveurs (+ système de\\ stockage)\end{tabular}     & Disponibilité       & \begin{tabular}[c]{@{}l@{}}- Dysfonctionnement de\\ matériel\\- Surtension électrique\\ - Personne mal-intentionnée\end{tabular}                                                                                                & \begin{tabular}[c]{@{}l@{}}- Perte des données\\ - Système inaccessible\\ - Réputation de l'hôpital\\ - Vie des patients\end{tabular} & Elevé    \\ \hline

	\begin{tabular}[c]{@{}l@{}}Dysfonctionnement du\\ système d'authentification\end{tabular} & Disponibilité       & \begin{tabular}[c]{@{}l@{}}- Erreur logiciel\\ - Utilisateur\end{tabular}                                                                                                                              & \begin{tabular}[c]{@{}l@{}}- Impossibilité de consulter\\ les informations\end{tabular}                                               & Elevé    \\ \hline

	\begin{tabular}[c]{@{}l@{}}Intrusion d'une personne\\ non autorisée\end{tabular}          & Confidentialité     & \begin{tabular}[c]{@{}l@{}}- Personnel qui a oublié\\ son badge et utilise celui\\ d'un collègue\\ - Mauvaise gestion des roles\\ assignés aux utilisateurs\\- Personne mal-intentionnée\end{tabular} & \begin{tabular}[c]{@{}l@{}}- Divulgation de données\\ personnelles à une personne\\ non-autorisée. (secret médical)\end{tabular}      & Elevé    \\ \hline

	\begin{tabular}[c]{@{}l@{}}Panne de courant\end{tabular}          & Disponibilité     & \begin{tabular}[c]{@{}l@{}}- Problème sur le réseau\\ électrique\\- Personne mal-intentionnée \end{tabular} & \begin{tabular}[c]{@{}l@{}}- Impossible de récupérer\\ les informations des patients \end{tabular}      & TODO    \\ \hline
	
	\begin{tabular}[c]{@{}l@{}}TODO\end{tabular}          & TODO     & \begin{tabular}[c]{@{}l@{}} TODO \end{tabular} & \begin{tabular}[c]{@{}l@{}}TODO \end{tabular}      & TODO    \\ \hline

\end{tabular}}
\caption{Table d'analyse de l'accès aux information médicales}
\label{table:accesInformationsMedicales}
\end{table}

\subsubsection{Encodage des données}

L'analyse de ce bien essentiel concerne l'encodage des données sur un patient. Par exemple, en fin de service par un membre du corps médical.

\begin{table}[ht]
\centering
\resizebox{\textwidth}{!}{%
\begin{tabular}{|l|l|l|l|l|}
\hline
Evenements redoutés & Critère de Sécurité & Source de la Menace & Impact & Severité \\ \hline
\begin{tabular}[c]{@{}l@{}}Intrusion d'une personne\\ non-autorisée\end{tabular} & Intégrité & \begin{tabular}[c]{@{}l@{}}- Personnel qui a oublié\\ son badge et utilise celui\\ d'un collègue\\ - Mauvaise gestion des roles\\ assignés aux utilisateurs\\ - Personne mal-intentionnée\end{tabular} & \begin{tabular}[c]{@{}l@{}}- Altération des données\\ des patients\\ - La vie du patient est\\ mise en danger\end{tabular} & Elevé \\ \hline
\begin{tabular}[c]{@{}l@{}}Dysfonctionnement du\\ système d'authentification\end{tabular} & Disponibilité & \begin{tabular}[c]{@{}l@{}}- Erreur logiciel\\ - Utilisateur\end{tabular} & - Système inutilisable & Elevé \\ \hline
Panne de courant & Disponibilité & \begin{tabular}[c]{@{}l@{}}- Problème sur le réseau électrique\\ - Personne mal-intentionnée\end{tabular} & - Système inutilisable & Elevé \\ \hline
\begin{tabular}[c]{@{}l@{}}Incendie dans la salle des\\ serveurs (+ système de\\ stockage)\end{tabular} & Disponibilité & \begin{tabular}[c]{@{}l@{}}- Dysfonctionnement du matériel\\ - Surtension électrique\\ - Personne mal-intentionnée\end{tabular} & - Système inutilisable & Elevé \\ \hline
 &  &  &  &  \\ \hline
\end{tabular}%
}
\caption{Table d'analyse de l'encodage des données}
\label{table:encodageDonnees}
\end{table}

\subsubsection{Réputation}

\subsection{Scénarios de menace}

Cette section est dédiée à une analyse des scénarios de menaces. Cette analyse est basée sur les biens support de l'hôpital.

\subsubsection{Serveurs}

\begin{table}[ht]
\resizebox{\textwidth}{!}{%
\begin{tabular}{|l|l|l|}
Scénario de menace & Source de la menace & Probabilité \\
Panne du serveur & \begin{tabular}[c]{@{}l@{}}Obsolescence du\\ matériel\end{tabular} & Moyenne \\
Incendie &  &  \\
 &  & 
\end{tabular}%
}
\end{table}

\subsubsection{Système d'authentification}

\subsubsection{Données}

\subsubsection{Réseau interne}

\subsubsection{Disques durs}

\subsubsection{Ordinateurs}

\subsection{Synthèse}

\section{Plan d'action}

\subsection{sub-1}

\subsection{sub-2}

\newpage

%---------------------------------------List of Figures----------------------------------------

\section*{Table des figures}
\makeatletter
\@starttoc{lof}% Print List of Figures
\makeatother
\addcontentsline{toc}{section}{\protect\numberline{}Table des figures}%

%---------------------------------------References----------------------------------------

\nocite{*}
\section*{Bibliographie}
\printbibliography[heading=none]
\addcontentsline{toc}{section}{\protect\numberline{}Bibliographie}%

%----------------------------------------------------------------------------------------
%----------------------------------------THE END-----------------------------------------
%----------------------------------------------------------------------------------------

\end{document}
